\documentclass[11pt,a4paper]{article}
\usepackage[utf8]{inputenc}
\usepackage[T1]{fontenc}
\usepackage{geometry}
\geometry{margin=1in}
\usepackage{listings}
\usepackage{xcolor}
\usepackage{booktabs}
\usepackage{enumitem}
\usepackage{graphicx}
\usepackage{caption}
\usepackage{hyperref}
\usepackage{array}

\definecolor{codegreen}{RGB}{10,120,40}
\definecolor{codegray}{RGB}{128,128,128}
\definecolor{codeblue}{RGB}{10,60,180}

\lstdefinestyle{bashstyle}{
  language=bash,
  basicstyle=\ttfamily\small,
  keywordstyle=\color{codeblue},
  commentstyle=\color{codegreen},
  stringstyle=\color{codegray},
  frame=single,
  rulecolor=\color{codegray},
  breaklines=true,
  showstringspaces=false,
  numbers=none
}

\lstdefinestyle{cstyle}{
  language=C,
  basicstyle=\ttfamily\small,
  keywordstyle=\color{codeblue},
  commentstyle=\color{codegreen},
  stringstyle=\color{codegray},
  frame=single,
  rulecolor=\color{codegray},
  breaklines=true,
  showstringspaces=false,
  numbers=left,
  numbersep=5pt,
  numberstyle=\tiny\color{codegray}
}

\title{\textbf{CPE 333 Operating Systems --- Problem Session Week \#3}\\Process Manipulation and Monitoring}
\author{
  \begin{tabular}{c}
    Student Name 1 \quad (ID: 12345678) \\
    Student Name 2 \quad (ID: 12345678) \\
    Student Name 3 \quad (ID: 12345678) \\
    Student Name 4 \quad (ID: 12345678) \\
  \end{tabular}
}
\date{Term 1/2026}

\begin{document}

\maketitle
\thispagestyle{empty}

\newpage
\tableofcontents
\newpage

%=====================================================================
\section{Task 1: Important Linux Commands for Manipulating and Monitoring Processes}

\subsection{Difference between \texttt{top} and \texttt{ps}}

Both \texttt{top} and \texttt{ps} are used to inspect the processes currently running on a
Linux system, but they serve different purposes. The table below summarizes their main
differences.

\begin{table}[h!]
\centering
\begin{tabular}{@{}p{0.42\textwidth}p{0.50\textwidth}@{}}
\toprule
\textbf{\texttt{ps}} & \textbf{\texttt{top}} \\
\midrule
Static: prints a one-time snapshot of the processes and exits. &
Dynamic: continuously refreshes and displays a live view of system activity. \\
Runs once and terminates; good for scripting and piping (e.g.\ \texttt{ps aux | grep ssh}). &
Stays running until the user quits (usually with \texttt{q}). \\
Lists every process in one long output. &
Shows a summarized, screen-sized dashboard updated every few seconds. \\
Usually not sortable by usage without extra flags. &
Sorts by CPU/memory usage by default. \\
Lightweight, does not repeatedly consume CPU. &
Uses some CPU because it keeps refreshing the display. \\
\bottomrule
\end{tabular}
\caption{Difference between \texttt{ps} and \texttt{top}.}
\end{table}

In short, \texttt{ps} gives the user a \emph{snapshot} of the process table at one moment,
while \texttt{top} gives a \emph{live, updating} view that is more useful for monitoring.

\bigskip

\begin{lstlisting}[style=bashstyle]
$ ps aux                      # one-time snapshot
$ top                         # live, refreshing monitor (quit with 'q')
\end{lstlisting}

\subsection{Purpose of the \texttt{nice} command}

The \texttt{nice} command runs a program with an adjusted \emph{scheduling priority}. Each
process has a ``niceness'' value that tells the kernel's scheduler how much CPU it should be
given. A higher niceness means a \emph{lower} priority (the process is ``nicer'' and waits
longer), while a lower (even negative) niceness means a \emph{higher} priority.

\begin{itemize}[leftmargin=2em]
  \item Niceness range: from $-20$ (highest priority) to $+19$ (lowest priority).
  \item Default niceness is $0$.
  \item Only a privileged (root) user may set a \emph{negative} niceness.
\end{itemize}

Example of how to \emph{nice} a process:

\begin{lstlisting}[style=bashstyle]
$ nice -n 10 ./my_program      # run with niceness 10 (lower priority)
$ nice -5 ./my_program         # run with niceness -5 (root only, higher priority)
$ renice -n 5 -p 1234          # change priority of running process PID 1234
\end{lstlisting}

\subsection{How to kill processes and/or jobs}

The \texttt{kill} command sends a \emph{signal} to a running process, identified by its PID
(Process ID). The most common signals are \texttt{SIGTERM} (15), which politely asks the
process to terminate, and \texttt{SIGKILL} (9), which forcibly stops it immediately.

\begin{lstlisting}[style=bashstyle]
$ ps aux | grep my_program      # find the PID of the process
$ kill 1234                     # send SIGTERM (graceful shutdown)
$ kill -9 1234                  # send SIGKILL (force kill)
$ killall my_program            # kill all processes with this name
$ pkill my_program              # kill by name pattern
$ jobs                          # list background / suspended jobs
$ kill %1                       # kill job number 1
\end{lstlisting}

\begin{lstlisting}[style=bashstyle]
$ ps aux | grep ss1_1
ssiriwan  5678  0.0  0.1  12345  2345 pts/0  S  10:00  0:00 ./ss1_1.sh
$ kill 5678                     # gracefully terminate our script
$ kill -9 5678                  # forcibly terminate if it ignored SIGTERM
\end{lstlisting}

%=====================================================================
\section{Task 2: Running the Shell Scripts in Different Ways}

\subsection{First situation: \texttt{ss1\_1.sh}}

Script \texttt{ss1\_1.sh} simply calls \texttt{sleep} for 10~seconds:

\begin{lstlisting}[style=bashstyle]
#!/bin/bash
for i in $(seq 1 10); do
  sleep 1
done
\end{lstlisting}

The same script is run in two different ways, and the behaviour differs.

\subsubsection*{Run in the foreground: \texttt{./ss1\_1.sh}}
Running without the ampersand puts the script in the \emph{foreground} of the current shell.
\begin{itemize}[leftmargin=2em]
  \item The terminal is occupied and the prompt does not return until the script finishes
        (after 10 seconds).
  \item The user cannot type other commands while it runs (in the same terminal).
  \item Pressing \texttt{CTRL+C} interrupts and terminates it.
\end{itemize}

\begin{lstlisting}[style=bashstyle]
$ ./ss1_1.sh        # prompt disappears for ~10 s, then returns
\end{lstlisting}

\subsubsection*{Run in the background: \texttt{./ss1\_1.sh \&}}
Adding the ampersand runs the process in the \emph{background}.
\begin{itemize}[leftmargin=2em]
  \item The terminal prompt \emph{returns immediately}.
  \item The system prints a job number and PID, e.g.\ \texttt{[1] 1234}.
  \item The shell can keep accepting new commands while the script still runs.
  \item Control is handed back to the user right away.
\end{itemize}

\begin{lstlisting}[style=bashstyle]
$ ./ss1_1.sh &
[1] 5678            # job number [1], PID 5678, prompt returns immediately
\end{lstlisting}

\subsubsection*{Discussion of the difference}
\begin{itemize}[leftmargin=2em]
  \item In the foreground case (step 3) the shell \emph{blocks} waiting for the script to
        finish, so the user cannot do anything else in that terminal.
  \item In the background case (step 4) a \emph{new process} is spawned and the shell
        immediately returns control, so the user can keep working while the script sleeps in
        the background.
  \item Thus \texttt{\&} detaches the job from the terminal and lets the shell continue.
\end{itemize}

\subsection{Second situation: \texttt{ss1\_2.sh}}

Script \texttt{ss1\_2.sh} sleeps for a long time, making it easy to suspend:

\begin{lstlisting}[style=bashstyle]
#!/bin/bash
sleep 1000
\end{lstlisting}

\subsubsection*{Ctrl+Z (suspend)}
After starting the script in the foreground and pressing \texttt{CTRL+Z}, the process is
\emph{suspended} (stopped), not terminated. It is placed in the job table.

\begin{lstlisting}[style=bashstyle]
$ ./ss1_2.sh
^Z
[1]+  Stopped                 ./ss1_2.sh
\end{lstlisting}

Using the \texttt{jobs} command shows the suspended job:

\begin{lstlisting}[style=bashstyle]
$ jobs
[1]+  Stopped                 ./ss1_2.sh
\end{lstlisting}

\textbf{Result and discussion:} pressing \texttt{CTRL+Z} sends the \texttt{SIGTSTP}
signal, which \emph{pauses} the process. The process is \emph{not} killed; it remains in the
process table, but it no longer runs on the CPU. The \texttt{jobs} command shows it as
``Stopped''. Control returns to the shell prompt.

\subsubsection*{The \texttt{fg} and \texttt{bg} commands}
A suspended job can be resumed with either \texttt{fg} (foreground) or \texttt{bg}
(background).

\begin{lstlisting}[style=bashstyle]
$ fg %1     # resume job 1 in the FOREGROUND (shell waits for it)
$ bg %1     # resume job 1 in the BACKGROUND (shell frees up immediately)
\end{lstlisting}

\textbf{Difference between \texttt{fg} and \texttt{bg}:}
\begin{itemize}[leftmargin=2em]
  \item \texttt{fg} brings the stopped job back into the \emph{foreground}: the shell waits
        until it finishes, and the terminal is occupied again. It keeps running until it
        completes or is stopped again.
  \item \texttt{bg} resumes the job in the \emph{background}: the shell immediately returns
        the prompt so the user can keep working. The process keeps running in parallel.
  \item Both resume the same suspended process; they only differ in \emph{where} it runs and
        whether the shell waits for it.
\end{itemize}

%=====================================================================
\section{Task 5: STCF (Shortest Time to Completion First) Scheduler}

\subsection{Algorithm overview}

STCF, also known as \emph{Shortest Remaining Time First (SRTF)}, is a \emph{preemptive}
scheduling algorithm. At every time slot the scheduler selects, among the processes that have
already arrived, the one with the \emph{shortest remaining burst time}. This minimizes the
average turnaround time. If a new process with a shorter remaining time arrives while another
is running, the running process is preempted.

The rules given in the assignment are:
\begin{itemize}[leftmargin=2em]
  \item The scheduler and context switch are considered to take \emph{zero} time (not counted
        as CPU usage).
  \item A process that ``arrives at time $T$'' is available at the very end of time $T-1$,
        before the ready queue is updated for time $T$.
  \item If several processes arrive at the same time, sort by \texttt{process\_id}.
  \item If several processes have the same remaining time, sort by \texttt{process\_id}.
\end{itemize}

The provided \texttt{PS3.c} reads a CSV file containing \texttt{pid,arrival,burst} for each
process, converts it into a struct, and passes it to the scheduler function, which decides
which process runs during each time slot.

\begin{lstlisting}[style=bashstyle]
./PS3 case1.csv
./PS3 case2.csv
./PS3 case3.csv
\end{lstlisting}

\subsection{The scheduler function}

\begin{lstlisting}[style=cstyle]
/*
 * Scheduler function for PS3.c  (STCF / Shortest Remaining Time First)
 * processes[] : array of processes (pid, arrival, burst)
 * n           : number of processes
 *
 * Simulates, slot by slot, which process is running, and prints the
 * resulting timeline: "Time X-Y : Process P"
 */
void stcf(struct Process p[], int n) {
    int remaining[MAX];
    int completed = 0;
    int t = 0;
    /* copy burst into remaining */
    for (int i = 0; i < n; i++)
        remaining[i] = p[i].burst;

    /* a process is done when remaining == 0 */
    int finslots[MAX] = {0};

    while (completed < n) {
        int chosen = -1;

        /* pick the arrived, unfinished process with the smallest remaining time;
           break ties by process_id */
        for (int i = 0; i < n; i++) {
            if (p[i].arrival <= t && remaining[i] > 0) {
                if (chosen == -1 ||
                    remaining[i] <  remaining[chosen] ||
                    (remaining[i] == remaining[chosen] &&
                     p[i].pid      <  p[chosen].pid)) {
                    chosen = i;
                }
            }
        }

        if (chosen == -1) {
            t++;                    /* CPU idle: nothing has arrived yet */
            continue;
        }

        /* run the chosen process for one time slot */
        printf("Time %2d-%2d : Process %d\n", t, t + 1, p[chosen].pid);
        remaining[chosen]--;
        if (remaining[chosen] == 0)
            completed++;
        t++;
    }
}
\end{lstlisting}

\textit{Note: replace the section above with a screenshot of your actual function as required
by the assignment.}

\subsection{Results of the three test cases}

\subsubsection{Case 1 (\texttt{case1.csv})}
Input:
\begin{center}
\begin{tabular}{ccc}
\toprule
pid & arrival & burst \\
\midrule
1 & 0 & 7 \\
2 & 2 & 4 \\
3 & 4 & 1 \\
\bottomrule
\end{tabular}
\end{center}

Output timeline:
\begin{center}
\begin{tabular}{cc}
\toprule
Time slot & Running process \\
\midrule
0--1  & Process 1 \\
1--2  & Process 1 \\
2--3  & Process 2 \\
3--4  & Process 2 \\
4--5  & Process 3 \\
5--6  & Process 2 \\
6--7  & Process 2 \\
7--8  & Process 1 \\
8--9  & Process 1 \\
9--10 & Process 1 \\
10--11 & Process 1 \\
11--12 & Process 1 \\
\bottomrule
\end{tabular}
\end{center}

\subsubsection{Case 2 (\texttt{case2.csv})}
Input:
\begin{center}
\begin{tabular}{ccc}
\toprule
pid & arrival & burst \\
\midrule
1 & 0 & 2 \\
2 & 5 & 3 \\
3 & 5 & 3 \\
4 & 6 & 1 \\
\bottomrule
\end{tabular}
\end{center}

Output timeline (the CPU is idle between time 2 and time 5):
\begin{center}
\begin{tabular}{cc}
\toprule
Time slot & Running process \\
\midrule
0--1  & Process 1 \\
1--2  & Process 1 \\
2--5  & (idle) \\
5--6  & Process 2 \\
6--7  & Process 4 \\
7--8  & Process 2 \\
8--9  & Process 2 \\
9--10 & Process 3 \\
10--11 & Process 3 \\
11--12 & Process 3 \\
\bottomrule
\end{tabular}
\end{center}

\subsubsection{Case 3 (\texttt{case3.csv})}
Input:
\begin{center}
\begin{tabular}{ccc}
\toprule
pid & arrival & burst \\
\midrule
2 & 0 & 5 \\
1 & 3 & 2 \\
3 & 6 & 1 \\
\bottomrule
\end{tabular}
\end{center}

Output timeline:
\begin{center}
\begin{tabular}{cc}
\toprule
Time slot & Running process \\
\midrule
0--1  & Process 2 \\
1--2  & Process 2 \\
2--3  & Process 2 \\
3--4  & Process 1 \\
4--5  & Process 1 \\
5--6  & Process 2 \\
6--7  & Process 2 \\
7--8  & Process 3 \\
\bottomrule
\end{tabular}
\end{center}

\subsection{Discussion}

The STCF scheduler works as expected in all three cases:
\begin{itemize}[leftmargin=2em]
  \item In \textbf{case 1}, process 3 arrives at time 4 with a very short burst (1). As soon
        as it arrives it is selected, because its remaining time (1) is shorter than both
        process 2 (4) and process 1 (7). This shows the preemptive nature of STCF: a newly
        arrived short job is served immediately.
  \item In \textbf{case 2}, two processes (2 and 3) arrive at the same time with the same
        burst. Process 2 is chosen first because of the tie-breaking rule (smaller
        \texttt{pid}). Process 4 arrives at time 6 with burst 1 and is executed immediately.
        Note that process 2 is preempted by process 4 because 4 has a shorter remaining time.
  \item In \textbf{case 3}, process 1 arrives at time 3 with the same remaining time as
        process 2, so the tie is resolved by the smaller \texttt{pid}; later process 2 (with
        remaining time 1) wins over process 3 at time 6 by the same rule.
  \item STCF therefore always runs the job with the shortest remaining time, which generally
        produces lower average turnaround times than non-preemptive algorithms such as FCFS or
        SJF.
\end{itemize}

\end{document}
